\subsection{The star oracle at scale}\label{sec:star-bench}

No block of Sections~\ref{sec:minlplib}--\ref{sec:mechanism} had more than
four variables, so the star algorithm of Theorem~\ref{thm:star} was not needed
there. Part~S measures it directly, without a solver, on random constrained
stars with $k\in\{10,100,1000,10000\}$ leaves (five instances each), $2k$
center--leaf rows, one to three rows in the center alone, and small rational
coefficients of both signs, so that convex, concave and affine leaves occur.
We compare an implementation of the sweep of the proof with the simpler star
oracle of our separator (Section~\ref{sec:star}). Where both ran ($k\le1000$),
the two minima were equal as rational numbers, and on all 20 instances the
sweep's minimizer was exactly feasible and attained the reported value. The
sweep took a median of 1.3\,ms, 12\,ms, 0.17\,s and 2.1\,s of CPU time for the
four values of $k$, close to linear growth, whereas the simpler oracle took
6\,ms, 0.43\,s and 29\,s for $k\le1000$, growing by a factor of about 70 per
tenfold increase in $k$; at $k=10^4$ it was not run, because its extrapolated
time exceeded ten minutes per instance. Its partitions had about twice as
many pieces as the sweep's (median 900 against 416 at $k=1000$). The minima
had numerators and denominators of at most 94 and 78 bits, with no steady
growth in $k$, in line with the size bounds of Appendix~\ref{app:star}.
